\subsection{\texorpdfstring{复形 \(\mathsf{K}(u),\mathsf{K}.(u,\mathsf{C}),\mathsf{K}^{\prime}(u,\mathsf{C})\)}{复形 \textbackslash mathsf\{K\}(u),\textbackslash mathsf\{K\}.(u,\textbackslash mathsf\{C\}),\textbackslash mathsf\{K\}\^{}\{\textbackslash prime\}(u,\textbackslash mathsf\{C\})}}

设 A 为环，L 为 A-模， \(u: \mathbf{L} \to \mathbf{A}\) 为线性形式，且 \(\Lambda(\mathbf{L})\) 为 A-模 L 的外代数。对于 \(x \in \Lambda(\mathbf{L})\) ，令 \(d_u(x)\) 表示内积 \(x \sqsubseteq u\) （III，第 161 页，例）。根据前引文第 162 页公式 (60)，我们有

\[
d _ {u} \left(e _ {1} \wedge \dots \wedge e _ {n}\right) = \sum_ {i = 1} ^ {n} (- 1) ^ {i + 1} u \left(e _ {i}\right) e _ {1} \wedge \dots \wedge e _ {i - 1} \wedge e _ {i + 1} \wedge \dots \wedge e _ {n}\tag{1}
\]

对 L 中的 \(e_1, \ldots, e_n\) 成立。根据 III，第 164 及 165 页，映射 \(d_u: \symbf{\Lambda}(\mathbf{L}) \to \symbf{\Lambda}(\mathbf{L})\) 是次数为 \((-1)\) 且平方为零的反导子。它是 A-代数 \(\symbf{\Lambda}(\mathbf{L})\) 中扩张 \(u: \symbf{\Lambda}^1(\mathbf{L}) \to \symbf{\Lambda}^0(\mathbf{L})\) 的唯一反导子。

定义 1.--- 复形 \((\symbf{\Lambda}(\mathbf{L}), d_u)\) 记作 \(\mathbf{K}^{\mathbf{A}}(u)\) 或 \(\mathbf{K}(u)\) 。

注意 \(\mathbf{K}_{n}(u)=\symbf{\Lambda}^{n}(\mathbf{L})=\mathbf{K}^{-n}(u)\) 。显然 \(\mathbf{K}(u)\) 右零，且 \(\mathrm{H}_{0}(\mathbf{K}(u))=\mathrm{Coker}(u)=\mathrm{A}/\mathfrak{q}\) ，其中 \(\mathfrak{q}\) 是 A 的理想 \(u(\mathbf{L})\) 。

对于 A-模的每个复形 C，我们设

\[
\mathbf {K} _ {\cdot} ^ {\mathrm{A}} (u, \mathrm{C}) = \mathrm{C} \otimes_ {\mathrm{A}} \mathbf {K} ^ {\mathrm{A}} (u), \quad \mathbf {K} _ {\cdot} ^ {\cdot} (u, \mathrm{C}) = \operatorname{Homgr} _ {\mathrm{A}} (\mathbf {K} ^ {\mathrm{A}} (u), \mathrm{C}).
\]

\[
\mathrm{H} ^ {\mathsf {A}} (u, \mathrm{C}) = \mathrm{H} (\mathrm{C} \otimes_ {\mathsf {A}} \mathbf {K} ^ {\mathsf {A}} (u)) \quad \mathrm{H} _ {\mathsf {A}} ^ {\bullet} (u, \mathrm{C}) = \mathrm{H} (\mathrm{Homgr} _ {\mathsf {A}} (\mathbf {K} ^ {\mathsf {A}} (u), \mathrm{C})) ,
\]

\[
\mathrm{H} _ {r} ^ {\mathrm{A}} (u, \mathrm{C}) = \mathrm{H} _ {r} (\mathrm{C} \otimes_ {\mathrm{A}} \mathbf {K} ^ {\mathrm{A}} (u)), \quad \mathrm{H} _ {\mathrm{A}} ^ {r} (u, \mathrm{C}) = \mathrm{H} ^ {r} (\operatorname{Homgr} _ {\mathrm{A}} (\mathbf {K} ^ {\mathrm{A}} (u), \mathrm{C})) .
\]

因此，我们有 A-模的典范同态（X，第 62 页及第 82 页）

\[
\gamma_ {0}: \mathrm{H} _ {0} (\mathbf {C}) \otimes_ {\mathbf {A}} \mathrm{A} / \mathfrak {q} \rightarrow \mathrm{H} _ {0} ^ {\mathbf {A}} (u, \mathbf {C}),
\]

\[
\lambda^ {0}: \mathrm{H} _ {\mathrm{A}} ^ {0} (u, \mathrm{C}) \rightarrow \operatorname{Hom} _ {\mathrm{A}} (\mathrm{A} / \mathfrak {q}, \mathrm{H} ^ {0} (\mathrm{C})).
\]

引理 1. --- 若复形 C 右零（相应地，左零），则 \(\mathbf{K}^{\mathbf{A}}(u,\mathbf{C})\) （相应地， \(\mathbf{K}_{\mathbf{A}}^{*}(u,c)\) ）右零（相应地，左零），并且 \(\gamma_{0}\) （相应地， \(\lambda^{0}\) ）是双射的。

这由 X，第 62 页，命题 1 及第 82 页，命题 1 得出。

命题 1. --- 设 \(x \in L\) ；记 \(\mathbf{R}_{x}: y \mapsto x \wedge y\) 为由 \(x\) 在代数 \(\boldsymbol{\Lambda}(\mathbf{L})\) 中左乘的映射。则有 \(d_{u} \circ \mathbf{R}_{x} + \mathbf{R}_{x} \circ d_{u} = u(x)\) . \(1_{\boldsymbol{\Lambda}(\mathbf{L})} = u(x)_{\boldsymbol{\Lambda}(\mathbf{L})}\) .

事实上， \((d_u \circ R_x + R_x \circ d_u)(y) = d_u(x \wedge y) + x \wedge d_u(y)\) ；由于 \(d_u\) 是一个导子， \(d_u(x \wedge y) + x \wedge d_u(y) = d_u(x) \wedge y = u(x).y\) .

推论 1. --- 若 u 是满射的，则 \(\mathbf{K}(u)\) 同伦于零（X，第 34 页）， \(\mathbf{K}_{\cdot}^{\mathrm{A}}(u,\mathrm{C})\) 和 \(\mathbf{K}_{\cdot}^{\cdot}(\mathbf{u},\mathbf{C})\) 对每个复形 C 也如此。

事实上，存在 \(x \in \mathbf{L}\) 使得 \(u(x) = 1\) 。于是 \(\mathbf{K}(u)\) 由命题 1 同伦于零，因而 \(\mathbf{K}^{\mathrm{A}}(u, \mathbf{C})\) （X，第 64 页，命题 3）与 \(\mathbf{K}_{\mathrm{A}}^{*}(u, \mathbf{C})\) （X，第 83 页，命题 3）也如此。

推论 2. --- 设 C 是一个复形，Ann(C) 为其零化子。则 q + Ann(C) 零化 H \(^{A}\) (u, C) 与 H \(^{*}_{A}\) (u, C)。

对每个 \(\lambda \in q\) ，位似 \(\lambda_{\mathbf{K}(u)}\) 由该命题同伦于零，因此 \(1_{C} \otimes \lambda_{\mathbf{K}(u)}\) 以及 Homgr \((\lambda_{\mathbf{K}(u)}, 1_{C})\) 依据 X，第 64 页，命题 3 及 X，第 83 页，命题 3 也同伦于零；由此可知 \(\lambda\) 零化 H.(u, C) 与 H'(u, C)。若 \(\lambda \in \text{Ann}(C)\) ，则 \(1_{\mathbf{K}(u)} \otimes \lambda_{C}\) 以及 Homgr \((1_{\mathbf{K}(u)}, \lambda_{C})\) 为零。

假设 L 是投射的（相应地， \(\mathbf{K}(u)\) 在次数 \(>0\) 时零调）。则复形 \(\symbf{\Lambda}(\mathbf{L})\) 是投射的（相应地，是 A/q 的一个分解），依据 III，第 87 页，推论 2；由 X，第 102 页（相应地，第 100 页），因此对每个 A-模 M，我们有同态

\begin{enumerate}
\def\labelenumi{(\arabic{enumi})}
\setcounter{enumi}{1}
\tightlist
\item
\end{enumerate}

\[
\mathrm{H} _ {r} ^ {\mathbf {A}} (u, \mathrm{M}) \rightarrow \operatorname{Tor} _ {r} ^ {\mathbf {A}} (\mathrm{A} / \mathrm{q}, \mathrm{M}), \quad \operatorname{Ext} _ {\mathbf {A}} ^ {r} (\mathrm{A} / \mathrm{q}, \mathrm{M}) \rightarrow \mathrm{H} _ {\mathbf {A}} ^ {r} (u, \mathrm{M})\tag{3}
\]

\[
\operatorname{Tor} _ {r} ^ {\mathrm{A}} (\mathrm{A} / \mathrm{q}, \mathrm{M}) \rightarrow \mathrm{H} _ {r} ^ {\mathrm{A}} (u, \mathrm{M}), \quad \mathrm{H} ^ {r} (u, \mathrm{M}) \rightarrow \operatorname{Ext} _ {\mathrm{A}} ^ {r} (\mathrm{A} / \mathrm{q}, \mathrm{M}).
\]

若 L 是投射的且 \(\mathbf{K}(u)\) 在次数 \(>0\) 时零调，则上述同态 (2) 和 (3) 是双射且互为逆（X，第 102 页，命题 1）。

命题 2. --- 设 \((\mathbf{L}_{i})_{i\in I}\) 为一族 A-模，其中集合 I 有限且全序。设 u 为 \(\bigoplus_{i\in I}\mathbf{L}_{i}\) 上的线性形式， \(u_{i}\) 为它在 \(\mathbf{L}_{i}\) 上的限制。

A-代数的典范同构（III，第 84 页）

\[
g: \bigotimes_{\substack{i\in I}}\symbf{\Lambda}(\mathrm{L}_{i})\to \symbf{\Lambda}(\bigoplus_{i\in I}\mathrm{L}_{i})
\]

是复形 \(\otimes \mathbf{K}(u_i)\) （X，第 63 页）到复形 \(\mathbf{K}(u)\) 上的同构。

事实上，根据 X，第 64 页，注记 4，复形 \(\otimes_{i\in I}\mathbf{K}(u_i)\) 的微分 D 是一个反导子；反导子 \(d_u\) 与 \(g\circ D\circ g^{-1}\) （ \(\boldsymbol{\Lambda}(\oplus L_i)\) 的）在 \(\oplus L_i\) 上与映射 \(x\mapsto u(x)\) . 1（从 \(\oplus L_i\) 到 \(\boldsymbol{\Lambda}(\oplus L_i)\) ）一致，因此它们相等（III，第 128 页，推论）。

设 C 和 C' 是两个 A-模复形。我们有（X，第 63 页及第 99 页）复形的典范同构

\[
\mathbf {C} \otimes_ {\mathbf {A}} \left(\mathbf {C} ^ {\prime} \otimes_ {\mathbf {A}} \mathbf {K} (u)\right)\rightarrow \left(\mathbf {C} \otimes_ {\mathbf {A}} \mathbf {C} ^ {\prime}\right) \otimes_ {\mathbf {A}} \mathbf {K} (u)
\]

\[
\operatorname{Homgr} _ {\mathbf {A}} \left(\mathrm{C} ^ {\prime}, \operatorname{Homgr} _ {\mathbf {A}} (\mathbf {K} (u), \mathrm{C})\right)\rightarrow \operatorname{Homgr} _ {\mathbf {A}} \left(\mathrm{C} ^ {\prime} \otimes_ {\mathbf {A}} \mathbf {K} (u), \mathrm{C}\right),
\]

即，同构

\begin{enumerate}
\def\labelenumi{(\arabic{enumi})}
\setcounter{enumi}{3}
\tightlist
\item
\end{enumerate}

\[
\mathbf {C} \otimes_ {\mathbf {A}} \mathbf {K} ^ {\mathrm{A}} (u, \mathbf {C} ^ {\prime}) \rightarrow \mathbf {K} ^ {\mathrm{A}} (u, \mathbf {C} \otimes_ {\mathbf {A}} \mathbf {C} ^ {\prime})\tag{5}
\]

\[
\operatorname{Homgr} _ {\mathbf {A}} \left(\mathrm{C} ^ {\prime}, \mathbb {K} _ {\mathbf {A}} ^ {*} (u, \mathrm{C})\right)\rightarrow \operatorname{Homgr} _ {\mathbf {A}} \left(\mathbb {K} _ {\mathbf {A}} ^ {*} (u, \mathrm{C} ^ {\prime}), \mathrm{C}\right).
\]

在 (4) 和 (5) 中，取 \(\mathbf{C}^{\prime} = \mathbf{K}(u^{\prime})\) ，其中 \(u^{\prime}: L^{\prime} \to A\) 是 A-模 \(L^{\prime}\) 上的线性形式，并注意到 \(\mathbf{K}^{\mathrm{A}}(u, \mathbf{K}(u^{\prime}))\) 按定义等于 \(\mathbf{K}(u^{\prime}) \otimes_{\mathrm{A}} \mathbf{K}(u)\) ，根据命题 2，它等同于 \(\mathbf{K}(u^{\prime} \oplus u)\) ，其中 \(u^{\prime} \oplus u: L^{\prime} \oplus L \to A\) 是线性形式 \((x^{\prime}, x) \mapsto u^{\prime}(x^{\prime}) + u(x)\) 。于是我们得到复形的同构

\begin{enumerate}
\def\labelenumi{(\arabic{enumi})}
\setcounter{enumi}{5}
\tightlist
\item
\end{enumerate}

\[
\mathbf {K} _ {\bullet} ^ {\mathrm{A}} (u ^ {\prime} \oplus u, \mathrm{C}) \rightarrow \mathbf {K} _ {\bullet} ^ {\mathrm{A}} (u, \mathbf {K} _ {\bullet} ^ {\mathrm{A}} (u ^ {\prime}, \mathrm{C}))\tag{7}
\]

\[
\mathbf {K} _ {\mathrm{A}} ^ {\bullet} (u ^ {\prime}, \mathbf {K} _ {\mathrm{A}} ^ {\bullet} (u, \mathrm{C})) \rightarrow \mathbf {K} _ {\mathrm{A}} ^ {\bullet} (u ^ {\prime} \oplus u, \mathrm{C}).
\]

通过取同调，我们推出A-模的同构

\[
\begin{array}{c c} \mathrm{H} _ {r} ^ {\mathbf {A}} (u ^ {\prime} \oplus u, \mathrm{C}) \to \mathrm{H} _ {r} ^ {\mathbf {A}} (u, \mathbf {K} _ {\bullet} ^ {\mathbf {A}} (u ^ {\prime}, \mathrm{C}))  , & r \in \mathbf {Z}  , \\ \mathrm{H} _ {\mathbf {A}} ^ {r} (u ^ {\prime}, \mathbf {K} _ {\mathbf {A}} ^ {\bullet} (u, \mathrm{C})) \to \mathrm{H} _ {\mathbf {A}} ^ {r} (u ^ {\prime} \oplus u, \mathrm{C})  , & r \in \mathbf {Z}  . \end{array}
\]

最后，注意到由代数 \(\Lambda(L)\) 中的乘法导出的同态

\[
m: \mathbf {K} ^ {\mathrm{A}} (u) \otimes_ {\mathrm{A}} \mathbf {K} ^ {\mathrm{A}} (u) \rightarrow \mathbf {K} ^ {\mathrm{A}} (u)
\]

是复形态射（因为 \(d_{u}\) 是反导子）。假设 L 是秩为 \(n\) 的自由模，并与复形态射 \(\mathbf{K}^{\mathrm{A}}(u)\to\symbf{\Lambda}^{n}\mathrm{L}(-n)\) （它在次数 \(n\) 处为恒同）复合，由此导出一个复形态射

\[
\chi : \mathbf {K} ^ {\mathrm{A}} (u) \otimes_ {\mathrm{A}} \mathbf {K} ^ {\mathrm{A}} (u) \rightarrow \boldsymbol {\Lambda} ^ {n} \mathrm{L} (- n);
\]

根据 X，第 99 页，命题 12，该态射典范地对应一个复形态射

\[
\varphi : \mathbf {K} ^ {\mathrm{A}} (u) \rightarrow \operatorname{Homgr} _ {\mathrm{A}} \left(\mathbf {K} ^ {\mathrm{A}} (u), \boldsymbol {\Lambda} ^ {n} \mathrm{L} (- n)\right)
\]

它是双射的（III，第 87 页，公式 (20)）。对每个复形 C，由此导出一个复合同构

\[
\begin{array}{r l} \mathbf {K} _ {\bullet} ^ {\mathrm{A}} (u, \mathrm{C}) & = \mathrm{C} \otimes_ {\mathrm{A}} \mathbf {K} ^ {\mathrm{A}} (u) \xrightarrow {1 \otimes_ {\Phi}} \mathrm{C} \otimes \operatorname{Homgr} _ {\mathrm{A}} \left(\mathbf {K} ^ {\mathrm{A}} (u), \boldsymbol {\Lambda} ^ {n} \mathrm{L} (- n)\right) \to \\ & \to \operatorname{Homgr} _ {\mathrm{A}} \left(\mathbf {K} ^ {\mathrm{A}} (u), \mathrm{C} \otimes_ {\mathrm{A}} \boldsymbol {\Lambda} ^ {n} \mathrm{L} (- n)\right) = \mathbf {K} _ {\mathrm{A}} ^ {\bullet} (u, \mathrm{C} \otimes_ {\mathrm{A}} \boldsymbol {\Lambda} ^ {n} \mathrm{L} (- n)). \end{array}
\]

通过取同调，因此有典范同构

\[
\mathrm{H} _ {r} ^ {\mathbf {A}} (u, \mathrm{C}) \rightarrow \mathrm{H} _ {\mathbf {A}} ^ {n - r} (u, \mathrm{C} \otimes_ {\mathbf {A}} \boldsymbol {\Lambda} ^ {n} \mathrm{L}), \quad r \in \mathbf {Z}.\tag{8}
\]

注记。--- 1) * 当 L 为秩 n 的投射模时，前述结论仍然成立。~*

\begin{enumerate}
\def\labelenumi{\arabic{enumi})}
\setcounter{enumi}{1}
\tightlist
\item
  由于 L 是秩为 n 的自由模， \(\Lambda^{n}\) L 同构于 A，故有非典范同构 \(\mathrm{H}_{r}^{\mathrm{A}}(u,\mathrm{C})\to\mathrm{H}_{\mathrm{A}}^{n-r}(u,\mathrm{C})\) 。
\end{enumerate}

